\documentclass[11pt]{article}
\usepackage{amsmath, amssymb, amsthm, enumerate}
\usepackage[margin=1in, top=1.2in, bottom=1.2in]{geometry}
\usepackage{parskip}
\usepackage[colorlinks=true, linkcolor=blue, citecolor=blue, urlcolor=blue]{hyperref}
\usepackage{cleveref}
\theoremstyle{definition}
\newtheorem{definition}{Definition}[section]
\newtheorem{theorem}{Theorem}[section]
\newtheorem{lemma}{Lemma}[section]
\newtheorem{proposition}{Proposition}[section]
\newtheorem{corollary}{Corollary}[section]
\theoremstyle{remark}
\newtheorem{remark}{Remark}[section]
\newtheorem{example}{Example}[section]
\setlength{\parindent}{0pt}
\setlength{\parskip}{6pt}
\begin{document}

\begin{center}
{\Large\textbf{Random Matrices}}\\
\vspace{0.3cm}
{\normalsize\today}
\end{center}
\vspace{0.5cm}


% --- Page 1 ---
---

\begin{definition}\label{def:random_matrix}
Consider a matrix of random entries:
\[
\begin{pmatrix}
* & \cdots & * \\
\vdots & \ddots & \vdots \\
* & \cdots & *
\end{pmatrix}
\]
These entries are random variables. Let \(X_1, \ldots, X_n\) be vectors in \(\mathbb{R}^d\), where each \(X_i = (X_{i1}, \ldots, X_{id})\) consists of random variables.
\end{definition}

\begin{definition}\label{def:variance_covariance}
For a matrix \(X\) with entries \(X_{ij}\):
\begin{enumerate}
    \item The \textbf{variance} of the \(i\)-th component is given by:
    \[
    \frac{1}{n} \sum_{i=1}^n X_{i,i}^2
    \]
    \item The \textbf{covariance} between components \(i_1\) and \(i_2\) is given by:
    \[
    \frac{1}{n} \sum_{i=1}^n X_{i,i_1} X_{i,i_2}
    \]
\end{enumerate}
\end{definition}

Let \(X\) be the data matrix:
\[
X = \begin{pmatrix}
X_{11} & \cdots & X_{1d} \\
\vdots & \ddots & \vdots \\
X_{n1} & \cdots & X_{nd}
\end{pmatrix}
\]
The sample covariance matrix is defined as:
\[
\frac{XX^\top}{n}
\]

\begin{definition}\label{def:eigenvalues}
The \textbf{eigenvalues} of a matrix \(A\) are defined by the equation:
\[
Av = \lambda v \quad \text{for} \quad v \neq 0.
\]
The trace of \(A\) is given by:
\[
\text{tr}(A) = \sum_i A_{ii} = \sum_i \lambda_i,
\]
where \(\lambda_i\) are the eigenvalues of \(A\).
\end{definition}

\begin{remark}\label{rem:spectral_distribution}
The set of eigenvalues of a random matrix is called the \textbf{empirical spectral distribution}. This is given by:
\[
\frac{1}{n} \sum_{i=1}^n \delta_{\lambda_i},
\]
where \(\delta_{\lambda_i}\) is the Dirac delta function centered at \(\lambda_i\).
\end{remark}

---

% --- Page 2 ---
The handwritten content in the image does not explicitly contain definitions, theorems, proofs, or other structured mathematical statements. Instead, it appears to be a set of notes on properties of a matrix \( A \) with entries that are probabilistic and symmetric.


\[
\begin{pmatrix}
A_{11} & A_{1j} & \cdots \\
A_{j1} & A_{jj} & \cdots \\
\vdots & \vdots & \ddots \\
A_{m1} & A_{mj} & \cdots & A_{mn}
\end{pmatrix}
\]

where the entries \( A_{ij} \) are probabilistic with \( \mathbb{P}(A_{ij} = 1) = \mathbb{P}(A_{ij} = -1) = \frac{1}{2} \) and \( A_{ji} = A_{ij} \).

\[
\text{tr}(A) = \sum_i A_{ii} = \sum_i \lambda_i
\]

\[
\text{tr}(A^2) = \sum_i A_{ii}^2 = \sum_i \lambda_i^2
\]

\[
\text{tr}(A^k) = \sum_i A_{ii}^k = \sum_i \lambda_i^k
\]

\end{document}